\documentclass[12pt,english]{article}
\usepackage[T1]{fontenc}
\usepackage[latin9]{inputenc}
\usepackage[toc,page]{appendix}
\usepackage{geometry}
\geometry{verbose,tmargin=1.5in,bmargin=1.5in,lmargin=1in,rmargin=1in}
\usepackage{bm}
\usepackage{amstext}
\usepackage{comment}
\usepackage{amsmath}
\usepackage{hyperref}
\hypersetup{
     colorlinks = true,
     citecolor = navy,
     linkcolor = navy,
     urlcolor=navy
}
\usepackage[space]{grffile}
\usepackage{setspace}
\usepackage[normalem]{ulem}
\usepackage{xcolor}
\usepackage{soul}
\usepackage{dsfont}
\usepackage{refstyle}
\usepackage[capitalize,nameinlink]{cleveref}
\usepackage{setspace}
\usepackage{bbm}
\usepackage[tableposition=t]{caption} % ensures caption of longtables is normalsized.
\usepackage{mathpazo} % add Palatino font
\usepackage{longtable}
\usepackage{amssymb}
\usepackage{amsfonts}
\usepackage{amsmath}
%\usepackage[margin=1in]{geometry} % to set margins
\usepackage{scalefnt}
\usepackage{wrapfig}
\usepackage{graphicx}
%\usepackage{subfigure}
\usepackage{lscape}
\usepackage{pdflscape}
\usepackage{multirow}
\usepackage{array}
\usepackage{rotating}
\usepackage{bibunits}
\usepackage{fix-cm}
\usepackage{float}
\floatplacement{table}{H}
\usepackage{epsfig}
\usepackage{epstopdf}
\usepackage{refstyle}
\usepackage{booktabs}
\usepackage{tabularx}
\usepackage{array}

\newcolumntype{x}[1]{>{\raggedleft\arraybackslash}p{#1}}%
\newcolumntype{d}[1]{>{\centering\arraybackslash}p{#1}}%
\newcolumntype{y}[1]{>{\raggedright\arraybackslash}p{#1}}%
\usepackage{verbatim}
%\epstopdfsetup{suffix=}
\setcounter{MaxMatrixCols}{10}
\usepackage[longnamesfirst,authoryear]{natbib}
\newenvironment{proof}[1][Proof]{\noindent \textbf{#1.}}{\rule{0.5em}{0.5em}}
\setlength{\topmargin}{-0.80in} \setlength{\textheight}{9.50in}
\setlength{\evensidemargin}{.0in} \setlength{\oddsidemargin}{.0in}
\setlength{\textwidth}{6.85in}
\setlength{\voffset}{0.4cm}
\setlength{\hoffset}{-0.5cm}
\usepackage{setspace}
\usepackage{afterpage}
\usepackage[flushleft]{threeparttable}
\usepackage{tabulary}
%\usepackage[titletoc,title]{appendix}
\usepackage[bottom]{footmisc}
\usepackage{color}
	\definecolor{maroon}{RGB}{128,0,0}
	\definecolor{navy}{RGB}{0,0,165}
	\definecolor{lb}{RGB}{150,165,190}
	\definecolor{MyWhite}{RGB}{255, 255, 255}
\usepackage{colortbl}

\usepackage{caption}
\usepackage{subcaption}
\usepackage{placeins}
\usepackage{nicefrac}
%\usepackage{chngcntr}

\newcommand{\bit}{\begin{itemize}}
\newcommand{\eit}{\end{itemize}}
\usepackage{listings}
\def\cN{\mathcal{N}}
\newtheorem{theorem}{Theorem}
\newtheorem{remark}[theorem]{Remark}
\usepackage{mdframed}
\newmdtheoremenv{assumption}{Assumption}

\usepackage{enumitem}


\begin{document}


%\onehalfspacing


\title{Running the VAR for ShapirIO}
%Beggar or Prosper-thy-Neighbour? \\
%Fiscal Spillovers Across Firms and Countries
\maketitle

\paragraph{Purpose:} The purpose of this note is to introduce some additional details for the ShapirIO project.

What we should currently have:

\begin{enumerate}
	\item Residuals for product x time for price and quantity: $\mu_{i,t}^{price},\;\mu_{i,t}^{quantity}$.
	\item An I-O matrix that is $I$ x $I$ where $I$ is the number of sectors. Let $IO_{i,j}$ be this matrix where $i$ {\bf \underline{buys from}} $j$. \textcolor{red}{Emily! This is different from before -- please transpose the other matrix!!!!!}
	\item A final demand vector that is $I$ x 1 and records final consumption. We have not saved final investment or government spending or exports. This might need to be fixed later on... Not sure! Let this vector be denoted as $\underline{FD}$
\end{enumerate}

What I would like you to do:

\begin{enumerate}
	\item Calculate an I-O matrix \emph{intermediate} cost shares:
	
	\[
		\Omega_{i,j} = \frac{IO_{i,j}}{\sum_{j} IO_{i,j}}
	\]
	Note that rows of this should sum to 1:
	\[
		\mathbf{\Omega}\; \underline{1}_{I}=\underline{1}_I
	\]
	where $1_I$ is a vector of ones and $\mathbf{\Omega}$ is the matrix of $\Omega_{i,j}$.
	\item Also calculate the intermediate share of sales:
	\[
		\gamma_{i}\equiv \frac{\sum_j IO_{i,j}}{\sum_j IO_{j,i}+FD_i}
	\]
	Note here that the numerator sums across a row and the denominator sums across a column of $IO$. In tensor product math this is:
	
	\[
		\underline{\gamma}\equiv (\mathbf{IO}\;\underline{1}_I)\oslash (\mathbf{IO}'\;\underline{1}_I+\underline{FD})
	\]
	where $\underline{x}$ means a vector of $x_i$ and $a \oslash b$ means that every element of $a$ is divided by the equivalent element of $b$.
	
	\item With this two concepts, calculate the following IO-adjusted price residual:
	\[
		\underline{u}^{price,IO}\equiv (I-diag(\underline{\gamma})\mathbf{\Omega})^{-1}\underline{u}^{price}
	\]
	\item Next let's solve for the IO-adjusted quantity shock. This takes two steps. First we calculate the change in sales for each sector (accounting for direct effects of changes in consumer spending and indirect effects through the I-O matrix). 
	\[
		\underline{u}^{sales,IO} = (I-\Omega'diag(\underline{\gamma}))^{-1}(e^{\underline{u}^{price}}\circ e^{\underline{u}^{quantity}})
	\]
	where $a\circ b$ means elementwise mutiplication of $a$ and $b$ ($.*$ in matlab, $*$ in Python).
	\item Next let's find the IO adjusted quantity residual:
	\[
		\underline{u}^{q,IO} = \log\left(\underline{u}^{sales,IO}\oslash e^{\underline{u}^{price,IO}}\right)
	\]
	where $a \oslash b$ denotes elementwise \emph{division} of $a$ in the numerator and $b$ in the denominator. This is equivalent to $./$ in Matlab and $/$ in Python.
\end{enumerate}

Then once you have these do the following:

\begin{itemize}
	\item Calculate the following: $u^{price,Difference}_{i,t} \equiv u_{i,t}^{price}-u_{i,t}^{price,IO}$. Then do a stacked bar graph where we decompose the Shapiro prices into the IO part and the difference. To make these look clear, can we do the following:
	\begin{itemize}
		\item Pick 1 month/quarter every decade and then do separate subplots for each decade (so there should be like 6 graphs). 
		\item On each graph, order the graphs by $u_{i,t}^{price}$. Then plot the stacked graphs.
		\item N.B. Stacked graph commands typically don't do negatives and positives together well. I suggest plotting all positive values first, then plot the negative values on top of them. When plotting the bars use $u_{i,t}^{price,IO}$ and $u_{i,t}^{difference,IO}$ only 
		\item Plot a line that corresponds to $u_{i,t}^{price}$ so we can see the ``total'' effect as a line. The final graph should look like Figure 4 at this link \href{https://www.rbnz.govt.nz/-/media/project/sites/rbnz/files/publications/analytical-notes/2013/an2013-09.pdf#page=13}{here} with only 2 colors.	
	\end{itemize}
	\item Repeat these graphs for the $u_{i,t}^{quantity}$ residuals 
	\item Do the Shapiro Classification: $I_{demand}=1$ if $sign(u_{i,t}^{price})=sign(u_{i,t}^{quantity})$ and $I_{supply}=1$ if $sign(u_{i,t}^{quantity})=-sign(u_{i,t}^{price})$
	\item Do the same classification for the IO residuals
	\item Do a heatmap where sectors are rows and there are two columns. Color column 1 green if there is a Shapiro demand shock and red if there is a Shapiro supply shock. Then color column 2 green if there is an IO demand shock and red if there is a IO supply shock.
	\item See if you can figure out a way to get the heatmap to also do negative/positive demand/supply shocks (4 colors perhaps)?
\end{itemize}

\section*{What to do when one NAICS has $n$ products?}

Here is my suggested approach:

\begin{enumerate}
	\item Start by creating a variable counting the number of products each NAICS sector is merged to.
	\item Create a variable that is called \emph{weight} that is 1 when the previous count variable is 1.
	\item Then groupby (NAICS) and create a variable that is the sum of all \emph{final expenditure} values with the same NAICS. If a NAICS code has one product then there will be a single expenditure entry (remember expenditure varies by product only not NAICS). 
	\item For cases where the count>1 (from step 1), set \emph{weight}=expenditure/sum(expenditure) where the sum is calculated from step 3.
	\item Set the IO value = IO value * weight 
	\item Now collapse the IO matrix to vary only by product
\end{enumerate}

\section*{Optimization Problem:}

\begin{align*}
	\max_{Y_1}&\;\;\; P_1Y_1 - w_1L_1 - P_1 M_{1\leftarrow 1}- P_2 M_{1\leftarrow 2}\\
		s.t.&\;\;\; Y_1=L_1^{\alpha}M_{1\leftarrow 1}^{\beta_1}M_{1\leftarrow 2}^{\beta_2}
\end{align*}

Let me know if you have any questions!


\end{document}